\section*{Closing Perspective}

This framework gives a decision-maker a disciplined way to locate where an AI-enabled workflow's risk actually lives --- in the model, in the tooling wrapped around it, in the organization's own habits of reliance, or in some combination that only becomes visible once the pieces are traced together. It gives a vocabulary precise enough to tell a prompt problem from a verification gap from a security-boundary failure, because those three are not fixed by the same intervention, and treating them as interchangeable is itself a common source of bad decisions. It gives a standard --- comparative reliability against a strong human baseline, not perfection --- specific enough to apply to an actual deployment rather than debate in the abstract.

What it cannot do is remove the judgment the decision actually requires. No framework establishes, on its own, what a specific firm's risk appetite should be, what a specific client relationship can bear, or when a specific residual uncertainty is small enough to accept. Those are exactly the calls this report is built to inform, not replace --- which is also why Section 15 and Appendix B close with genuinely open questions rather than manufactured resolution. An unresolved question is not, by itself, a reason to decline a deployment; a framework that claimed to have none would be less trustworthy than one that names them, not more. The right response to an open question is to ask whether it actually bears on the decision in front of you --- and if it doesn't, to proceed with a clear account of what remains unknown, rather than waiting for a certainty this document, or any other, cannot supply.
